\begin{itemize}
  \item \textbf{Reconocimiento confiable:} Las pruebas realizadas con el prototipo mostraron un alto nivel de precisión en la identificación de residuos plásticos y de cartón bajo condiciones controladas.
  \item \textbf{Interacción efectiva:} El uso de LEDs para indicar la clasificación correcta o incorrecta brindó retroalimentación clara y rápida a los usuarios, mejorando su experiencia.
  \item \textbf{Monitoreo en tiempo real:} El sistema de monitoreo del nivel de llenado demostró ser efectivo para notificar cuando el contenedor requiere ser vaciado, optimizando el manejo de desechos.
  \item \textbf{Adaptabilidad del diseño:} La estructura física y la programación del prototipo permiten integrarse en diferentes contextos, desde instituciones educativas hasta espacios comerciales.
  \item \textbf{Base para futuros desarrollos:} Los datos recolectados durante la operación del prototipo proporcionan una base sólida para realizar ajustes en el algoritmo de clasificación y expandir el sistema a nuevas funcionalidades.
\end{itemize}